% ============================================================
%   §7 Discussion
% ============================================================
\section{Discussion}
\label{sec:discussion}

The five Parts of the project pursue different questions through different methods --- graph topology, community detection, supervised learning, sparse precision estimation, cross-similarity analysis --- but several themes recur across them. This section pulls them together and discusses what they say about the dataset, the methods, and the limits of what can be concluded from the analysis.

\subsection*{The spatial graph is a productive abstraction at every scale}

The most consistent finding across the project is that the spatial KNN graph is a productive level of abstraction at every scale of analysis. \partref{1} extracts cohort-level structure from it (the tissue connectivity gradient, composition--topology correlations, betweenness skeletons). \partref{2} uses it as the substrate for Leiden community detection, which delivers 20 stable niches with biologically interpretable composition. \partref{3} feeds it to a GraphSAGE network and gains 6.5 percentage points of validation macro-$F_1$ over a strong per-cell baseline. Parts \ref{sec:part4}--\ref{sec:part5} fit gene networks on the niches that the graph induces, recovering biologically meaningful hub programmes.

No analysis required moving away from the KNN-with-radius-cap construction. This is not a trivial observation: in spatial transcriptomics there is a continuous design choice between Delaunay triangulation, KNN, fixed-radius graphs, learned graphs, and various hybrid constructions, each with different topological properties. The simple KNN-with-radius-cap used here proved sufficient for every analysis in the project. The implication is that the binding constraint on what can be learned from this kind of data is not graph construction; it is what the data themselves contain.

\subsection*{T cells are anomalous on every measure --- because they are spatially dispersed}

If a single result deserves to be flagged as cross-validated by independent analyses, it is the structural distinctness of T cells (\classid{6}). Each Part identifies it as different in a different way: \partref{1} flags it through negative cell-type assortativity (T cells sit in neighbourhoods dominated by other classes); \partref{2} observes that T cells never dominate a niche and appear as a substantial secondary minority in only \nicheid{17}, the B-cell-dominant niche likely corresponding to tertiary lymphoid structures; \partref{3} reports T cells as the worst-performing class across every rung of the ablation, with peak $F_1 \approx 0.42$ versus $\geq 0.88$ for the best class; \partref{4} finds that the T-cell gene network shares essentially zero structure with the seven other classes (top-15 hub Jaccard $\leq 0.03$, edge Jaccard $\leq 0.02$). The Graphical Lasso also fails to converge on T cells within 200 iterations.

The unifying explanation is biological. T cells in tumour tissue are infiltrating populations: they appear scattered in low density throughout multiple compartments (tumour, stroma, immune-rich regions), with no consistent spatial partner cell type. This spatial heterogeneity propagates into every downstream analysis. In topology terms (\partref{1}), it places T cells in neighbourhoods dominated by other cell types. In niche terms (\partref{2}), it prevents T cells from accumulating into a compositionally homogeneous niche of their own. In classification terms (\partref{3}), it makes T cells the hardest class to predict from neighbourhood features, because their neighbourhood is unstable. In gene-network terms (\partref{4}), it produces a microenvironmental gene programme that averages over many distinct neighbourhoods and lacks a coherent dominant signal --- and that is therefore disjoint from every other cell type's network. The Graphical Lasso convergence failure is the statistical signature of this heterogeneity: the covariance structure of T cells does not stabilise on a single dominant pattern within the iteration budget.

Notably, T cells emerge in only one niche as a substantial secondary population: \nicheid{17}, where they coexist with B cells at 29\% to 48\% proportions. This composition is consistent with tertiary lymphoid structures, organised B/T-cell aggregates that form in tumour tissue as part of the anti-tumour immune response. The fact that this is the \emph{only} place where T cells form a non-trivial population component, and that the same niche is the least edge-integrated of the B-cell-dominant niches in \partref{5} (a weak signal on a flat dendrogram, not a clean separation), is a coherent biological story: T cells in this dataset are diffuse infiltrators with one organised exception (the TLS-like niche), and the network-based analysis points --- weakly --- in the same direction.

\subsection*{Shared hubs, rewired connections at multiple scales}

A second cross-cutting finding is that the "shared hubs, rewired connections" pattern of gene networks emerges at both the cell-type level (\partref{4}) and the niche level (\partref{5}). At both levels, top-degree hub genes overlap moderately between similar groups (max hub Jaccard $0.36$ for classes, $0.58$ for niches), while the specific edges between those hubs overlap much less (max edge Jaccard $0.09$ for classes, $0.12$ for niches). The shared hubs across classes (CXCL8, S100A8/A9, COL1A2 among the universals) are best interpreted, in the diffusion-contamination framing of \partref{4}, as gene programmes that are broadly distributed across the tumour microenvironment: they appear as hubs in any class whose neighbourhood includes inflammatory cells or extracellular matrix. The "rewired connections" finding is the genuine class-specific signal; the "shared hubs" finding is partly a property of the data (broadly active programmes) and partly an artefact of mRNA diffusion across cell boundaries at near single-cell resolution.

A subsidiary finding from \partref{5} sharpens this: \emph{pairwise} network similarity correlates moderately with compositional similarity ($\rho = 0.615$), but \emph{cluster-level} structure recovery is poor (ARI $0.25$ for edges, $0.05$ for hubs). The networks contain enough signal to flag pairs of similar niches but not enough to reproduce the compositional family partition. This is a useful methodological caution: a Spearman correlation of $0.6$ between two distance matrices is not the same as agreement at the partition level, and confusing the two is a common error in cross-modality analyses.

\subsection*{The connectivity gradient imposes downstream constraints}

The variation in tissue connectivity documented in \partref{1} (largest connected component fraction spanning $0.11$ to $0.999$, in a continuous gradient rather than two distinct modes) has practical consequences throughout the project. Niche detection in \partref{2} works well on compact samples (median $96\%$ cell coverage) but only partially on fragmented samples (median $32\%$), because it operates on the giant component and the LCF therefore caps the achievable coverage. The GNN of \partref{3} is trained on a \texttt{train\_mask} that is itself density-biased (Spearman $\rho = -0.91$ between mask retention and mean degree). Per-niche GGMs in \partref{5} subsample from the training pool, which under-represents fragmented samples.

These constraints are not failures --- the analyses are still internally valid --- but they bound the cohort-level generalisation. A reader of this report should take away that the headline numbers (test macro-$F_1 = 0.6089$, $20$ stable niches, edge Jaccard $\rho = 0.615$ vs.\ composition) are valid summaries of the cohort as analysed, but they will tend to overrepresent compact-architecture tissue. Fragmented samples are present in the test set, but their internal structure is less precisely characterised by the niche framework.

\subsection*{The graph subsumes the H\&E embedding for cell typing}

A more specific methodological observation from \partref{3} is that the UNI H\&E embedding contributes $\sim 8.5$pp of validation $F_1$ inside a multi-modal MLP, but contributes nothing when added on top of GraphSAGE. The two modalities therefore overlap in the information they carry about cell-type identity. This is not surprising biologically: the H\&E embedding captures local tissue architecture (cell density, packing, adjacency to other cells), which is precisely what the spatial graph propagates through neighbour aggregation. Whichever of the two modalities is present first saturates the architectural signal; the second adds little.

For practitioners, this suggests that on cell-typing tasks where a spatial graph is available, expensive foundation-model embeddings of the H\&E image may be a poor return on compute. Conversely, in datasets where the spatial graph is unavailable or unreliable, the H\&E embedding partially substitutes. The two are alternative entry points to roughly the same architectural information.

\subsection*{Diffusion contamination shapes what the gene networks can tell us}

A property of the data that becomes visible only in Parts \ref{sec:part4}--\ref{sec:part5} but applies to every per-cell analysis in the project is \emph{mRNA diffusion contamination}: at the near single-cell resolution of Visium HD, mRNA from a target cell and its $\sim 1$--$2$ nearest spatial neighbours is difficult to disambiguate, and the expression vector of a cell is partly a vector of its spatial neighbourhood. This was confirmed externally by inspection of canonical marker gene dotplots, which show the expected diagonal pattern (each cell type expressing its canonical markers) overlaid on substantial off-diagonal signal from neighbouring cell types.

The implication for \partref{4} is that the per-class gene networks recover not pure cell-intrinsic regulatory programmes but \emph{microenvironmental gene programmes} --- a mix of the cell's own regulation and the gene expression of cells it typically sits next to. The strength of this mixing varies by class: Myeloid cells, Fibroblasts, SMC, Tumor, and Normal epithelium retain some canonical markers among their hubs (their cell-intrinsic signal survives because they cluster spatially with similar populations); B cells, Endothelial cells, and T cells are dominated by neighbourhood signal because they sit in heterogeneous spatial contexts. The biological interpretations developed in \partref{4} (B cells inheriting SMC signal from perivascular TLS contexts; Tumor cells inheriting CAF signal from the surrounding stroma; T cells unable to commit to any consistent neighbourhood) are consistent with this picture and with the spatial biology of tumour microenvironments.

This is not a methodological flaw of the project but a property of the platform. Spatial transcriptomics at near single-cell resolution will always carry some neighbourhood contamination, and the right way to read the per-class and per-niche GGMs is as conditional-dependence networks of \emph{neighbourhoods}, not of \emph{cell types in isolation}. The "shared hubs, rewired connections" finding is robust to this reading: pairs of cell types that share spatial neighbourhoods share gene network hubs, but the specific conditional-dependence edges remain class-specific. A future iteration of this analysis on the same platform would benefit from deconvolution methods (Bayesian or matrix-factorisation-based) that attempt to separate cell-intrinsic from neighbourhood signal at the read level; in the current report, the diffusion is acknowledged and incorporated into the interpretation rather than corrected.

\subsection*{What I cannot claim}

A few qualifications are worth stating explicitly.

First, cell-type hub interpretations. The interpretations of hub gene programmes in \partref{4} are inferences about what the GGM is recovering --- a mixture of cell-intrinsic regulation and diffusion-mediated neighbourhood signal --- and should be read as biologically plausible but not validated. In particular, the interpretation of \nicheid{17} as a tertiary lymphoid structure follows from its composition (B cells dominant, T cells substantial secondary, low frequency in the cohort) but is not confirmed by independent immunohistochemistry or transcript-level markers of TLS organisation.

Second, the niche count $K = 20$. This was chosen by inspection of the dendrogram and intra-cluster homogeneity. A nearby choice ($K = 15$ or $K = 25$) would merge or split some of the smaller niches but is unlikely to change the qualitative findings. I did not perform a formal stability analysis on $K$.

Third, the Graphical Lasso convergence failures. Endothelial cells, Fibroblasts, and T cells did not reach the dual-gap tolerance within 200 iterations for the per-class fits, and several per-niche fits had the same issue. The results from non-converged iterates appear stable in bootstrap, but I cannot rule out that a more careful regularisation schedule (or a different sparse-precision estimator) would yield different edge sets. The qualitative findings (hub identity, hub vs.\ edge Jaccard, class--niche concordance) are robust across the bootstrap resamples, which is the strongest stability evidence I have.

Fourth, the test-validation gap. Validation $\macroF = 0.6350$ vs.\ test $\macroF = 0.6089$ ($-2.6$pp gap) is small but non-zero. With only 8 test patients, the per-patient variance is non-trivial ($\sigma = 0.042$, range $[0.473, 0.625]$). A larger test set would tighten the headline estimate, but the current configuration is the standard one for the dataset.

Fifth, the Gaussianity assumption of the Graphical Lasso. Log-normalised spatial transcriptomics counts are zero-inflated and non-Gaussian; restricting to broadly variable common genes mitigates the issue at the marginal level but does not eliminate it. The principled alternative is the nonparanormal SKEPTIC estimator \cite{liu2012nonparanormal}, which uses rank correlations and is consistent under arbitrary continuous marginal transformations. A SKEPTIC refit would be the natural next step for any extension that derives quantitative conclusions from the precision-matrix entries; the qualitative findings of \partref{4}--\partref{5} (hub identity, shared-hubs vs.\ rewired-connections, T-cell outlier) are likely robust but not formally validated against the more defensible estimator.

\subsection*{Directions the analysis points toward}

A few questions are left open. The poor per-niche performance on niche-pure regions (\nicheid{0}, \nicheid{1}, \nicheid{15} all below median test $F_1$) suggests that a niche-conditioned cell-typing approach could help: instead of one classifier for all niches, train a niche-specialised head per family. \partref{5} already provides the per-niche partition needed for this. The differential gene-network analysis between paired similar niches (e.g.\ \nicheid{17} vs.\ \nicheid{16}/\nicheid{18}/\nicheid{19}, all B-cell-dominant) is also a natural extension: the differential edges between two compositionally similar niches localise the niche-specific regulatory rewiring that \partref{5} only captures at the aggregate level. A complementary direction is to compare the gene network of the B-cell-only \nicheid{16} against the B-cell-plus-T-cell \nicheid{17}: the network differences would isolate the conditional-dependence signature of the T-cell co-population, with potential implications for TLS organisation in tumour tissue. Both directions are well-motivated by the findings of the current analysis and would build directly on the artefacts already produced.

\medskip

\noindent\textbf{Closing remark.} The project demonstrates that a single dataset of $\sim 8$ million spatially-resolved cells supports a wide range of network-theoretic analyses: graph topology, community detection, neighbourhood-aware classification, and conditional-dependence networks of genes. The most informative result is not any individual headline number; it is the consistency with which the same patterns (T cells as the spatially dispersed outlier, shared-hubs rewired-connections, the tissue connectivity gradient constraining every downstream analysis, the dominance of the graph over per-cell modalities) reappear across independent analyses. The spatial graph is the integrating object, the conditional-independence networks are the most refined level at which biological structure becomes visible, and the diffusion contamination characteristic of high-resolution spatial transcriptomics shapes those networks into descriptions of \emph{microenvironments} rather than of cell types in isolation.